\chapter{Locally convex direct sums and direct limits}

Locally convex direct sums and direct limits, after \cite[Appendix~B.13 and Sections~2.7--2.8]{GN26}.
They are the tool for spaces of compactly supported functions and sections, fine box products of
manifolds and manifolds of mappings on non-compact manifolds \cite[Sections~4.3--4.4]{GN26}.

\section{Direct sums}

\begin{definition}[Locally convex direct sum]
  \label{def:ds-direct-sum}
  \lean{LocallyConvex.Analysis.LCDirectSum, LocallyConvex.Analysis.LCDirectSum.box,
    LocallyConvex.Analysis.LCDirectSum.addGroupFilterBasis}
  \leanok
  For topological vector spaces $(E_j)_{j \in J}$, $\bigoplus_j E_j$ is the space of finitely
  supported families with the \emph{box topology}: the boxes
  $\bigoplus_j V_j = \{x \mid x_j \in V_j \text{ for all } j\}$, $V_j$ a $0$-neighbourhood in $E_j$,
  form a basis of $0$-neighbourhoods.
\end{definition}

\begin{theorem}[Basic properties {\cite[Remark~B.13.1, Lemma~B.13.2(a), (c)]{GN26}}]
  \label{thm:ds-basic}
  \lean{LocallyConvex.Analysis.LCDirectSum.nhds_zero_hasBasis,
    LocallyConvex.Analysis.LCDirectSum.nhds_hasBasis,
    LocallyConvex.Analysis.LCDirectSum.isOpen_box,
    LocallyConvex.Analysis.LCDirectSum.continuous_apply,
    LocallyConvex.Analysis.LCDirectSum.continuous_single,
    LocallyConvex.Analysis.LCDirectSum.continuous_of_continuous_single,
    LocallyConvex.Analysis.LCDirectSum.prodEquiv,
    LocallyConvex.Analysis.LCDirectSum.sigmaEquiv}
  \leanok
  \uses{def:ds-direct-sum}
  The box topology is a vector topology over any nontrivially normed field, locally convex and
  Hausdorff when the $E_j$ are; boxes of open sets are open; coordinates and inclusions are
  continuous. For countable $J$, an additive map from $\bigoplus_j E_j$ to a topological group is
  continuous as soon as its restrictions to the summands are; hence the box topology is the locally
  convex direct sum topology of \cite[(B.24)]{GN26}. Moreover
  $\bigoplus_j (E_j \times F_j) \cong \bigoplus_j E_j \times \bigoplus_j F_j$ and, for any index
  sets, $\bigoplus_{(i, j)} E_{ij} \cong \bigoplus_i \bigoplus_j E_{ij}$
  \cite[Lemma~B.13.2(b)]{GN26}.
\end{theorem}

\begin{proof}
  \leanok
  For the universal property, as in \cite[Lemma~B.13.2(a)]{GN26}: choose $0$-neighbourhoods $Q_k$
  with $Q_{k+1} + Q_{k+1} \subseteq Q_k$ and an injection $n \colon J \to \mathbb{N}$; the box of
  preimages of $Q_{n_j + 1}$ is mapped into $Q_0$, by induction on the support, adding the summand of
  least index first. For uncountable $J$ the locally convex direct sum topology can be strictly finer
  than the box topology.
\end{proof}

\begin{theorem}[Subspaces, bounded sets, completeness {\cite[Lemma~B.13.2(d)--(f)]{GN26}}]
  \label{thm:ds-subspaces}
  \lean{LocallyConvex.Analysis.LCDirectSum.isEmbedding_mapCLM_subtype,
    LocallyConvex.Analysis.LCDirectSum.exists_finset_of_isVonNBounded,
    LocallyConvex.Analysis.LCDirectSum.tendsto_mk_of_finset,
    LocallyConvex.Analysis.LCDirectSum.seqComplete,
    LocallyConvex.Analysis.LCDirectSum.mackeyComplete,
    LocallyConvex.Analysis.LCDirectSum.quasiCompleteSpace}
  \leanok
  \uses{thm:ds-basic}
  If $F_j \subseteq E_j$ are subspaces, $\bigoplus_j F_j$ carries the topology induced by
  $\bigoplus_j E_j$. If the $E_j$ are Hausdorff, every bounded subset of $\bigoplus_j E_j$ lies in
  $\bigoplus_{j \in J_0} E_j$ for a finite $J_0 \subseteq J$, and $\bigoplus_j E_j$ is complete
  (quasi-complete, sequentially complete, Mackey complete) if every $E_j$ is.
\end{theorem}

\begin{proof}
  \leanok
  The box of the sets $V_j \cap F_j$ is the trace of the box of the $V_j$. If a bounded set met
  infinitely many summands, at coordinates $x_{j_k} \neq 0$, a box omitting $k^{-1} x_{j_k}$ would
  not absorb it. For completeness, the coordinates of a Cauchy filter converge, and the limit has
  finite support by the same argument. Closed bounded sets, Cauchy sequences and Mackey--Cauchy
  sequences are bounded, hence lie in a finite partial sum, where convergence is coordinatewise;
  the coordinates converge in the factors.
\end{proof}

\begin{theorem}[Direct sums of $\Ck{k}$ and real analytic maps {\cite[Proposition~2.8.1]{GN26}}]
  \label{thm:ds-calculus}
  \lean{LocallyConvex.Analysis.LCDirectSum.map,
    LocallyConvex.Analysis.LCDirectSum.continuousOn_map,
    LocallyConvex.Analysis.LCDirectSum.hasDirDerivAt_map,
    LocallyConvex.Analysis.LCDirectSum.contDiffMBOn_map,
    LocallyConvex.Analysis.LCDirectSum.smoothMBOn_map,
    LocallyConvex.Analysis.LCDirectSum.complexAnalyticOn_map,
    LocallyConvex.Analysis.LCDirectSum.realAnalyticOn_map}
  \leanok
  \uses{thm:ds-basic}
  Let $f_j \colon U_j \to F_j$ be $\Ck{k}$ ($k \in \mathbb{N} \cup \{\infty\}$) on open sets
  $U_j \subseteq E_j$, with $f_j(0) = 0$ for all but finitely many $j$. Then
  $\bigoplus_j f_j \colon (x_j)_j \mapsto (f_j(x_j))_j$ is $\Ck{k}$ on the open set
  $\bigoplus_j U_j$, with $\dd(\bigoplus_j f_j) = (\bigoplus_j \dd f_j) \circ \Phi$. If each $f_j$
  is real analytic and $0 \in U_j$ for all but finitely many $j$, then $\bigoplus_j f_j$ is real
  analytic.
\end{theorem}

\begin{proof}
  \leanok
  Induction on $k$, as in \cite{GN26}: the difference quotients of $\bigoplus_j f_j$ only involve
  finitely many coordinates. In the real analytic case $\bigoplus_j \tilde f_j$, for complex analytic
  extensions $\tilde f_j$, is smooth with complex linear differentials, hence complex analytic, and
  extends $\bigoplus_j f_j$ through $(\bigoplus_j E_j)_{\mathbb{C}} \cong \bigoplus_j (E_j)_{\mathbb{C}}$.
\end{proof}

\section{Direct limits}

\begin{theorem}[Locally convex direct limits {\cite[Definition~B.13.4, Remark~B.13.5]{GN26}}]
  \label{thm:ds-limit}
  \lean{LocallyConvex.Analysis.LCDirectLimit,
    LocallyConvex.Analysis.LCDirectLimit.continuous_incl,
    LocallyConvex.Analysis.LCDirectLimit.continuous_of_continuous_comp_incl,
    LocallyConvex.Analysis.LCDirectLimit.t2Space_of_injective}
  \leanok
  Let $S_n$ be subspaces of a real vector space $V$, each with a vector topology. The seminorms of
  $V$ whose restrictions to every $S_n$ are continuous define the locally convex direct limit
  topology. The inclusions are continuous, a linear map to a locally convex space is continuous as
  soon as its restrictions to the $S_n$ are, and an injective such map to a Hausdorff space makes the
  limit Hausdorff.
\end{theorem}

\begin{proof}
  \leanok
  For the universal property, compose with the gauge seminorms of the target.
\end{proof}

\begin{remark}[Cofinal sequences {\cite[Remark~B.13.10]{GN26}}]
  \label{rem:ds-cofinal}
  \lean{LocallyConvex.Analysis.LCDirectLimit.topology_eq_comp}
  \leanok
  \uses{thm:ds-limit}
  If a monotone family $(S_i)_{i \in I}$ of subspaces with continuous inclusions admits a cofinal
  sequence $(S_{i_n})_n$, then both have the same locally convex direct limit topology: apply the
  universal property in both directions.
\end{remark}

\begin{lemma}[Convex neighbourhoods {\cite[Lemma~B.13.6]{GN26}}]
  \label{lem:ds-convex-nhds}
  \lean{LocallyConvex.Analysis.LCDirectLimit.mem_nhds_zero_iff_of_convex,
    LocallyConvex.Analysis.LCDirectLimit.isOpen_iff_of_convex,
    LocallyConvex.Analysis.LCDirectLimit.barrelledSpace,
    LocallyConvex.Analysis.locallyConvexSpace_quotient,
    LocallyConvex.Analysis.LCDirectLimit.isOpenQuotientMap_sumMap,
    LocallyConvex.Analysis.LCDirectLimit.nhds_zero_hasBasis_sumMap,
    LocallyConvex.Analysis.LCDirectLimit.isOpen_sumMap_box}
  \leanok
  \uses{thm:ds-limit, thm:ds-basic}
  Let the $S_n$ cover $V$.
  \begin{enumerate}
    \item A convex set $U \ni 0$ is a $0$-neighbourhood of the direct limit if and only if
      $U \cap S_n$ is a $0$-neighbourhood of $S_n$ for every $n$.
    \item For an ascending sequence with continuous inclusions, a convex set is open if and only
      if its traces on the $S_n$ are open.
    \item For countably many steps, the summation map $\bigoplus_n S_n \to V$ is an open quotient
      map; hence the sets $\bigcup_n (W_0 + \cdots + W_n)$, $W_n$ (open) $0$-neighbourhoods of
      $S_n$, form a basis of (open) $0$-neighbourhoods.
    \item If the $S_n$ are barrelled, so is the direct limit.
  \end{enumerate}
\end{lemma}

\begin{proof}
  \leanok
  (a) The gauge of $U \cap (-U)$ is a seminorm on $V$ whose restrictions to the $S_n$ are
  continuous. (b) Translate to a point of $U$ and apply (a). (c) The summation map is continuous
  by the universal property of the direct sum; on the quotient of the direct sum by its kernel,
  a locally convex space, the inclusions of the $S_n$ are continuous, so the inverse of the
  induced bijection is continuous by the universal property of the direct limit. (d) A lower
  semicontinuous seminorm on the limit restricts to continuous seminorms on the steps.
\end{proof}

\begin{lemma}[Extension of absolutely convex open sets {\cite[Lemma~B.13.8]{GN26}}]
  \label{lem:ds-extension}
  \lean{LocallyConvex.Analysis.exists_absConvex_isOpen_preimage_eq}
  \leanok
  Let $j \colon F \to E$ be an injective linear topological embedding into a locally convex space.
  Every absolutely convex open set $V \subseteq F$ is $j^{-1}(Q)$ for an absolutely convex open set
  $Q \subseteq E$.
\end{lemma}

\begin{proof}
  \leanok
  As in \cite{GN26}, $Q$ is the convex join of $j(V)$ and an absolutely convex open
  $0$-neighbourhood $W$ of $E$ with $j^{-1}(W) \subseteq V$. A point $a\, j(v) + b w$ of $Q$ has the
  open neighbourhood $a\, j(v) + b W \subseteq Q$ if $b > 0$, and $j(v) = r^{-1} j(rv) + (1 - r^{-1})
  \cdot 0$ with $rv \in V$, $r > 1$, otherwise.
\end{proof}

\begin{theorem}[Strict direct limits {\cite[Proposition~B.13.9(a)--(d)]{GN26}}]
  \label{thm:ds-strict}
  \lean{LocallyConvex.Analysis.LCDirectLimit.IsStrict,
    LocallyConvex.Analysis.LCDirectLimit.IsStrict.exists_mem_nhds_preimage_eq_of_subset,
    LocallyConvex.Analysis.LCDirectLimit.IsStrict.isEmbedding_incl,
    LocallyConvex.Analysis.LCDirectLimit.IsStrict.t2Space,
    LocallyConvex.Analysis.LCDirectLimit.IsStrict.exists_subset_isVonNBounded,
    LocallyConvex.Analysis.LCDirectLimit.IsStrict.exists_subset_isCompact,
    LocallyConvex.Analysis.LCDirectLimit.IsStrict.completeSpace,
    LocallyConvex.Analysis.LCDirectLimit.IsStrict.quasiCompleteSpace,
    LocallyConvex.Analysis.LCDirectLimit.IsStrict.exists_tendsto_of_cauchySeq,
    LocallyConvex.Analysis.LCDirectLimit.IsStrict.mackeyComplete}
  \leanok
  \uses{lem:ds-convex-nhds, lem:ds-extension}
  Let $S_0 \subseteq S_1 \subseteq \cdots$ be locally convex subspaces covering $V$, each $S_n$
  carrying the topology induced by $S_{n+1}$. Then:
  \begin{enumerate}
    \item the direct limit induces the given topology on each $S_n$;
    \item the direct limit is Hausdorff if the $S_n$ are;
    \item if each $S_n$ is closed in $S_{n+1}$, every bounded (compact) subset of the limit is
      contained, and bounded (compact), in some $S_n$;
    \item the direct limit is complete if the $S_n$ are;
    \item if each $S_n$ is closed in $S_{n+1}$, the direct limit is quasi-complete (sequentially
      complete, Mackey complete) if the $S_n$ are.
  \end{enumerate}
\end{theorem}

\begin{proof}
  \leanok
  (a) Extend an absolutely convex open $0$-neighbourhood of $S_m$ step by step with
  Lemma~\ref{lem:ds-extension} and take the union, a $0$-neighbourhood by
  Lemma~\ref{lem:ds-convex-nhds}. (b) follows from (a). (c) If a bounded set $B$ met infinitely
  many differences $S_{n+1} \setminus S_n$, at points $x_n$, the extensions in (a) can be cut down
  so that $x_n \notin (n+1) W$, and $W$ does not absorb $B$. (d) As in \cite{GN26}, a Cauchy filter
  concentrates near some $S_m$: otherwise the convex hull of sets $W_n \cap S_n$, with $W_n$
  shrinking $0$-neighbourhoods, contradicts the Cauchy property. The traces on $S_m$ of the sets
  $A + W$ then form a Cauchy filter on $S_m$, by (a), whose limit is the limit of the filter.
  (e) Closed bounded sets, Cauchy sequences and Mackey--Cauchy sequences (together with their
  bounded set) are bounded, hence lie and are bounded in some $S_n$, by (c); conclude with (a).
\end{proof}

\section{Direct limits of topological spaces}

\begin{definition}[Direct sequences {\cite[Definitions~A.7.1, A.7.5, Remark~A.7.15]{GN26}}]
  \label{def:top-dl}
  \lean{LocallyConvex.Analysis.TopDirectLimit, LocallyConvex.Analysis.IsDirectSeq,
    LocallyConvex.Analysis.IsDirectSeq.IsStrict, LocallyConvex.Analysis.IsDirectSeq.IsCompactifying,
    LocallyConvex.Analysis.IsDirectSeq.CompactlyRegular}
  \leanok
  Let $X_0 \subseteq X_1 \subseteq \cdots$ be topological spaces covering a set $X$, with continuous
  inclusions. The \emph{direct limit topology} on $X$ is the final topology of the inclusions. The
  sequence is \emph{strict} if each $X_n$ carries the topology induced by $X_{n+1}$, and
  \emph{compactifying} if every point of $X_n$ has a neighbourhood whose closure in some $X_m$ is
  compact. It is \emph{compactly regular} if every compact subset of the limit lies, and is
  compact, in some $X_n$.
\end{definition}

\begin{theorem}[Basic properties {\cite[Remark~A.7.2, Lemmas~A.7.3, A.7.6, Remark~A.7.15]{GN26}}]
  \label{thm:top-dl-basic}
  \lean{LocallyConvex.Analysis.TopDirectLimit.isOpen_iff,
    LocallyConvex.Analysis.TopDirectLimit.continuous_iff,
    LocallyConvex.Analysis.TopDirectLimit.eq_topology_iff,
    LocallyConvex.Analysis.TopDirectLimit.topology_eq_comp,
    LocallyConvex.Analysis.TopDirectLimit.continuous_map,
    LocallyConvex.Analysis.IsDirectSeq.isOpen_iUnion,
    LocallyConvex.Analysis.IsDirectSeq.exists_subset_of_isCompact,
    LocallyConvex.Analysis.IsDirectSeq.IsStrict.isEmbedding_incl,
    LocallyConvex.Analysis.IsDirectSeq.IsStrict.compactlyRegular,
    LocallyConvex.Analysis.IsDirectSeq.isCompactifying_of_locallyCompactSpace,
    LocallyConvex.Analysis.IsDirectSeq.isStrict_of_compactSpace,
    LocallyConvex.Analysis.TopDirectLimit.topology_subtype_eq,
    LocallyConvex.Analysis.IsDirectSeq.topology_iUnion_eq}
  \leanok
  \uses{def:top-dl}
  A set is open (closed) in the direct limit iff its traces on all steps are; a map on the limit is
  continuous iff its restrictions to the steps are, and this characterizes the topology; cofinal
  subsequences give the same limit; ascending unions of open subsets of the steps are open. If the
  steps are $T_1$, so is the limit, and compact sets lie in some step. Strict sequences have
  embedded steps and, for $T_1$ steps, are compactly regular. Locally compact Hausdorff steps are
  compactifying; compact Hausdorff steps are strict. Open and closed subsets carry the direct limit
  topology of their traces, and so do ascending unions $\bigcup_n V_n$ of open subsets
  $V_n \subseteq X_n$, with respect to the $V_n$.
\end{theorem}

\begin{proof}
  \leanok
  As in \cite{GN26}. For compact sets, points $x_n \in K \setminus X_n$ would form an infinite
  closed discrete subset of $K$. For strict sequences, open subsets of $X_m$ are extended
  recursively to open subsets of the later steps.
\end{proof}

\begin{theorem}[Compactifying sequences {\cite[Lemmas~A.7.4, A.7.6(h), A.7.9,
  Proposition~A.7.20]{GN26}}]
  \label{thm:top-dl-compactifying}
  \lean{LocallyConvex.Analysis.IsDirectSeq.IsCompactifying.exists_isOpen,
    LocallyConvex.Analysis.IsDirectSeq.IsCompactifying.exists_disjoint,
    LocallyConvex.Analysis.IsDirectSeq.IsCompactifying.t2Space,
    LocallyConvex.Analysis.IsDirectSeq.IsCompactifying.isOpen_prod_iff,
    LocallyConvex.Analysis.IsDirectSeq.IsCompactifying.continuous_prod_iff,
    LocallyConvex.Analysis.IsDirectSeq.IsCompactifying.compactlyRegular,
    LocallyConvex.Analysis.IsDirectSeq.IsCompactifying.isLocallyKOmegaSpace,
    LocallyConvex.Analysis.isKOmegaSpace_iff_exists_topology_eq,
    LocallyConvex.Analysis.IsKOmegaSpace.isLocallyKOmegaSpace,
    LocallyConvex.Analysis.isKOmegaSpace_of_sigmaCompactSpace,
    LocallyConvex.Analysis.isLocallyKOmegaSpace_of_locallyCompactSpace,
    LocallyConvex.Analysis.IsLocallyKOmegaSpace.compactlyGeneratedSpace}
  \leanok
  \uses{thm:top-dl-basic}
  Let $X_0 \subseteq X_1 \subseteq \cdots$ and $Y_0 \subseteq Y_1 \subseteq \cdots$ be
  compactifying sequences of Hausdorff spaces. Then $\varinjlim X_n$ is Hausdorff, compactly
  regular and locally $k_\omega$, and $\varinjlim X_n \times \varinjlim Y_n =
  \varinjlim (X_n \times Y_n)$ as topological spaces. A Hausdorff space is $k_\omega$ iff it is
  the direct limit of an ascending sequence of compact subspaces; $k_\omega$-spaces and locally
  compact spaces are locally $k_\omega$, and locally $k_\omega$-spaces are $k$-spaces
  \cite[Remark~A.7.19]{GN26}.
\end{theorem}

\begin{proof}
  \leanok
  Compact subsets of a step have open neighbourhoods with compact closure in a later step, inside
  any prescribed open set. Iterating gives chains $K_j \subseteq V_j \subseteq K_{j+1}$, with $V_j$
  open and $K_j$ compact in increasing steps, whose unions are open: two disjoint chains separate
  points, and products of chains fit inside open subsets of the product. A set meeting every
  $K_j$ in a closed set is closed in $\bigcup_j V_j$, so compact subsets of $\bigcup_j V_j$ lie in
  some $K_j$ (otherwise they contain an infinite closed discrete set); covering a compact set by
  finitely many compact pieces of such unions gives compact regularity.
\end{proof}

\begin{theorem}[Groups and vector spaces {\cite[Proposition~A.7.11, Corollary~A.7.12,
  Example~A.7.13]{GN26}}]
  \label{thm:top-dl-groups}
  \lean{LocallyConvex.Analysis.TopDirectLimit.continuous_inv',
    LocallyConvex.Analysis.TopDirectLimit.continuous_mul_left',
    LocallyConvex.Analysis.TopDirectLimit.continuous_mul_right',
    LocallyConvex.Analysis.IsDirectSeq.IsCompactifying.isTopologicalGroup,
    LocallyConvex.Analysis.IsDirectSeq.IsCompactifying.nhds_one_hasBasis,
    LocallyConvex.Analysis.IsDirectSeq.isTopologicalGroup_of_locallyCompactSpace,
    LocallyConvex.Analysis.IsDirectSeq.IsCompactifying.continuousSMul,
    LocallyConvex.Analysis.IsDirectSeq.IsCompactifying.locallyConvexSpace,
    LocallyConvex.Analysis.IsDirectSeq.IsCompactifying.topology_eq_lcDirectLimit}
  \leanok
  \uses{thm:top-dl-compactifying, thm:ds-limit}
  For an ascending sequence of topological groups with continuous inclusions, inversion and
  translations are continuous in the direct limit topology. If the sequence is compactifying and
  the steps are Hausdorff, the limit is a topological group, with the sets
  $\bigcup_n U_0 U_1 \cdots U_n$ ($U_n$ open identity neighbourhoods of the steps) as a basis of
  identity neighbourhoods. Compactifying direct limits of Hausdorff (locally convex) topological
  vector spaces are (locally convex) topological vector spaces, and the direct limit topology is
  then the locally convex direct limit topology.
\end{theorem}

\begin{proof}
  \leanok
  Multiplication and scalar multiplication are continuous by the product theorem, the field being
  the constant compactifying sequence. In the basis the $U_n$ can be taken convex, and ascending
  unions of convex sets are convex. A locally convex topology making the inclusions continuous is
  coarser than the locally convex direct limit topology, by its universal property.
\end{proof}

\begin{theorem}[Further properties {\cite[Remark~A.7.7, Example~A.7.17, Remark~A.7.21,
  Lemma~A.7.22]{GN26}, \cite[Section~4]{GGH10}}]
  \label{thm:top-dl-komega}
  \lean{LocallyConvex.Analysis.IsDirectSeq.exists_nested_nhds,
    LocallyConvex.Analysis.IsDirectSeq.IsCompactifying.isCompactifying_unionSet,
    LocallyConvex.Analysis.compactlyGeneratedSpace_iff_topology_eq,
    LocallyConvex.Analysis.IsDirectSeq.t2Space_isKOmegaSpace,
    LocallyConvex.Analysis.IsDirectSeq.isLocallyKOmegaSpace,
    LocallyConvex.Analysis.IsKOmegaSpace.prod,
    LocallyConvex.Analysis.IsLocallyKOmegaSpace.prod,
    LocallyConvex.Analysis.IsDirectSeq.isOpen_prod_iff,
    LocallyConvex.Analysis.IsDirectSeq.isOpen_prod_iff_of_isKOmegaSpace}
  \leanok
  \uses{thm:top-dl-compactifying}
  For locally compact Hausdorff steps, a point of an open set has nested open neighbourhoods
  $V_n \subseteq X_n$ with compact closures inside the set and $\overline{V_n} \subseteq V_{n+1}$.
  The sets $V_n$ of Lemma~A.7.6(b) form a compactifying sequence if the $X_n$ do. A Hausdorff
  space is a $k$-space iff it is the direct limit of its compact subsets. Direct limits of
  (locally) $k_\omega$-spaces are (locally) $k_\omega$ and compactly regular, finite products of
  (locally) $k_\omega$-spaces are (locally) $k_\omega$, and
  $\varinjlim X_n \times \varinjlim Y_n = \varinjlim (X_n \times Y_n)$ for locally $k_\omega$
  steps.
\end{theorem}

\begin{proof}
  \leanok
  As in \cite{GN26} and \cite{GGH10}. For $k_\omega$ steps, the diagonal of cumulative
  $k_\omega$-sequences is an ascending sequence of compact spaces with the same direct limit. For
  locally $k_\omega$ steps, open $k_\omega$ sets $W_n \subseteq X_n$ with $W_n \subseteq W_{n+1}$
  (countable unions of open $k_\omega$ sets are $k_\omega$) have union $\varinjlim W_n$. The
  product of two such limits is a $k$-space, and its compact subsets lie in products of compact
  subsets of the steps.
\end{proof}

\begin{theorem}[The finite topology {\cite[Examples~A.7.14, A.7.16]{GN26}}]
  \label{thm:top-dl-finite}
  \lean{LocallyConvex.Analysis.FiniteTop, LocallyConvex.Analysis.FiniteTop.continuous_linearMap,
    LocallyConvex.Analysis.FiniteTop.topology_le,
    LocallyConvex.Analysis.FiniteTop.topology_eq_seq,
    LocallyConvex.Analysis.FiniteTop.locallyConvexSpace,
    LocallyConvex.Analysis.FiniteTop.continuousLinearEquivFinsupp}
  \leanok
  \uses{thm:top-dl-groups}
  The direct limit topology of the finite-dimensional subspaces of a real vector space $E$ is finer
  than every vector topology, and linear maps on it are continuous. If $E$ has countable
  dimension, it is a Hausdorff locally convex vector topology, hence the finest vector topology
  and the finest locally convex topology, and $E \cong \mathbb{R}^{(\mathbb{N})}$.
\end{theorem}

\begin{proof}
  \leanok
  Linear maps on finite-dimensional spaces are continuous. For countable dimension, the spans of
  the first $n$ basis vectors are cofinal and locally compact, so the limit is a locally convex
  vector topology by Corollary~A.7.12.
\end{proof}


\section{Silva spaces}

\begin{definition}[Silva systems {\cite[Definition~B.13.12]{GN26}}]
  \label{def:ds-silva}
  \lean{LocallyConvex.Analysis.IsSilvaSystem, LocallyConvex.Analysis.LCDirectLimit.IsSilva,
    LocallyConvex.Analysis.LCDirectLimit.IsSilva.isSilvaSystem}
  \leanok
  \uses{lem:ds-convex-nhds}
  A \emph{Silva system} presenting a locally convex space $E$ consists of continuous linear maps
  $j_n \colon E_n \to E$ from locally convex spaces and compact operators
  $u_n \colon E_n \to E_{n+1}$ with $j_{n+1} \circ u_n = j_n$, such that the ranges of the $j_n$
  cover $E$ and a convex set $C \ni 0$ is a $0$-neighbourhood as soon as every $j_n^{-1}(C)$ is.
  Silva spaces, locally convex direct limits of Banach spaces with compact inclusions, are the case
  of injective $j_n$ (by Lemma~\ref{lem:ds-convex-nhds}).
\end{definition}

\begin{theorem}[Topology of Silva spaces {\cite[Proposition~B.13.13(a)--(d)]{GN26}}]
  \label{thm:ds-silva}
  \lean{LocallyConvex.Analysis.IsSilvaSystem.isOpen_iff,
    LocallyConvex.Analysis.IsSilvaSystem.prod,
    LocallyConvex.Analysis.LCDirectLimit.IsSilva.isOpen_iff,
    LocallyConvex.Analysis.LCDirectLimit.IsSilva.continuous_iff,
    LocallyConvex.Analysis.LCDirectLimit.IsSilva.t2Space,
    LocallyConvex.Analysis.LCDirectLimit.IsSilva.isClosed_of_isClosed_inter,
    LocallyConvex.Analysis.LCDirectLimit.IsSilva.exists_subset_isCompact,
    LocallyConvex.Analysis.LCDirectLimit.IsSilva.exists_subset_isVonNBounded,
    LocallyConvex.Analysis.LCDirectLimit.IsSilva.isCompact_closure_of_isVonNBounded}
  \leanok
  \uses{def:ds-silva}
  Let $E$ be presented by a Silva system.
  \begin{enumerate}
    \item A set $O \subseteq E$ is open as soon as every $j_n^{-1}(O)$ is open: the locally convex
      direct limit topology is the direct limit topology of the topological spaces $E_n$. Hence a
      map from $E$ to a topological space is continuous as soon as its composites with the $j_n$
      are, and $E$ is Hausdorff if the $E_n$ are $T_1$ and the $j_n$ injective.
    \item Products of Silva systems are Silva systems.
  \end{enumerate}
  For a Silva space with $T_1$ sequential steps (for instance Banach spaces), compact subsets are
  contained, and compact, in some step; if moreover the steps have bounded $0$-neighbourhoods,
  bounded subsets are contained, and bounded, in some step, and they are relatively compact.
\end{theorem}

\begin{proof}
  \leanok
  We do not follow the route of \cite{GN26} through direct limits of compactifying sequences and
  $k_\omega$-spaces (Corollary~A.7.12, Proposition~A.7.20). (a) Given $x = j_m(y_0) \in O$, build
  compact sets $K_i \subseteq E_{m+i}$ with $j_{m+i}(K_i) \subseteq O$ and convex sets $A_i \ni 0$
  with $x + A_i \subseteq j_{m+i}(K_i)$: if $U_i$ is a convex $0$-neighbourhood of $E_{m+i}$ with
  $u_{m+i}(U_i)$ inside a compact set $L_i$, take $A_{i+1} = A_i + \varepsilon_i j_{m+i}(U_i)$ and
  $K_{i+1} = u_{m+i}(K_i) + \varepsilon_i L_i$, with $\varepsilon_i > 0$ so small that $K_{i+1}$
  is still mapped into $O$, by compactness. The union of the $A_i$ is a convex
  $0$-neighbourhood, and $x$ plus it lies in $O$. (b) A convex set $C \ni 0$ contains
  $\tfrac12 C_1 \times \tfrac12 C_2$, where $C_1 \times \{0\}$ and $\{0\} \times C_2$ are its
  traces on the factors. For compact regularity, the compact sets $k L_p$ cover $E$ and, by (a),
  detect closed sets; a compact set meeting the complements of all their finite unions would
  contain an infinite closed discrete set. Bounded regularity and relative compactness then follow
  as in \cite{GN26}, with an increasing sequence of bounded $0$-neighbourhoods of the steps in
  place of normalized unit balls. We do not claim these arguments are new.
\end{proof}

\begin{theorem}[Duality for Silva spaces {\cite[Proposition~B.13.13(e)--(g)]{GN26}}]
  \label{thm:ds-silva-duality}
  \lean{LocallyConvex.Analysis.exists_evalStrongDual_eq,
    LocallyConvex.Analysis.exists_continuousLinearEquiv_strongBidual,
    LocallyConvex.Analysis.LCDirectLimit.IsSilva.isKOmegaSpace,
    LocallyConvex.Analysis.LCDirectLimit.IsSilva.metrizableSpace_strongDual,
    LocallyConvex.Analysis.LCDirectLimit.IsSilva.completeSpace_strongDual,
    LocallyConvex.Analysis.LCDirectLimit.IsSilva.exists_continuousLinearEquiv_strongBidual,
    LocallyConvex.Analysis.LCDirectLimit.IsSilva.completeSpace}
  \leanok
  \uses{thm:ds-silva}
  A barrelled Hausdorff locally convex space whose bounded sets are relatively compact is
  reflexive. Consequently, a Silva space with $T_1$ sequential barrelled steps having bounded
  $0$-neighbourhoods (for instance Banach spaces) is a reflexive, complete $k_\omega$-space, and its
  strong dual is a Fréchet space.
\end{theorem}

\begin{proof}
  \leanok
  A functional $\varphi$ on the strong dual is bounded by $1$ on the functionals bounded by $1$ on
  a compact absolutely convex set $C$. If no point of $C$ matched $\varphi$ on finitely many
  functionals $f_1, \dots, f_k$, a functional on $\mathbb{R}^k$ would separate
  $(\varphi(f_i))_i$ from the image of $C$; compactness of $C$ then gives a point matching
  $\varphi$ everywhere. The evaluation map is continuous by Banach--Steinhaus and open by the
  bipolar argument (Hahn--Banach). For Silva spaces, the compact sets $k L_p$ detect closed sets
  and absorb the compact sets, so $E$ is $k_\omega$; its topology is coherent with the bounded
  sets, so the strong dual is complete, and a sequence of compact sets absorbing the bounded sets
  makes it metrizable. The strong bidual of a metrizable space is complete, so $E$ is complete.
  We do not claim this argument is new.
\end{proof}

\begin{theorem}[$\Ck{k}$ maps on Silva spaces {\cite[Proposition~2.7.1]{GN26}}]
  \label{thm:ds-silva-calculus}
  \lean{LocallyConvex.Analysis.IsSilvaSystem.contDiffMBOn_of_comp,
    LocallyConvex.Analysis.LCDirectLimit.IsSilva.contDiffMBOn_of_comp,
    LocallyConvex.Analysis.LCDirectLimit.IsSilva.contDiffMBOn_iUnion}
  \leanok
  \uses{thm:ds-silva}
  Let $E$ be presented by a Silva system with Hausdorff steps and $U \subseteq E$ open. A map
  $f \colon U \to F$ to a Hausdorff locally convex space is $\Ck{k}$ as soon as every
  $f \circ j_n$ is $\Ck{k}$ on $j_n^{-1}(U)$. In a Silva space, the union of ascending open subsets
  $U_n$ of the steps is open, and a map on it is $\Ck{k}$ as soon as its restrictions to the $U_n$
  are.
\end{theorem}

\begin{proof}
  \leanok
  Induction on $k$, as in \cite{GN26}: the directional derivatives of $f$ are those of the
  $f \circ j_n$, and $\dd f$ is $\Ck{k-1}$ on $U \times E$, presented by the product Silva system,
  by the induction hypothesis. Continuity comes from Theorem~\ref{thm:ds-silva}(a).
\end{proof}
